\documentclass[11pt,a4paper]{article}
\usepackage[T1]{fontenc}
\usepackage[utf8]{inputenc}
\usepackage{amsmath,amssymb,graphicx,hyperref,longtable,booktabs,array}
\usepackage[margin=2.2cm]{geometry}
\providecommand{\tightlist}{\setlength{\itemsep}{0pt}\setlength{\parskip}{0pt}}
\setlength{\parindent}{0pt}\setlength{\parskip}{6pt}

\begin{document}
\hypertarget{piece-2-realistic-sensor-layer-and-mismatch-generators-30-oct-2026}{%
\section{Piece 2: realistic sensor layer and mismatch generators (30 Oct
2026)}\label{piece-2-realistic-sensor-layer-and-mismatch-generators-30-oct-2026}}

Code: \texttt{code/sensor.py}, \texttt{code/test\_sensor.py},
\texttt{code/scenario\_report.py}. Input: the plasma traces in
\texttt{data/sim.npz} (1-min, 27 usable virtual patients). Written by
Claude (AI assistant); the team must read and sign off before claiming
it.

\hypertarget{what-it-does}{%
\subsection{What it does}\label{what-it-does}}

Builds a CGM-like signal from true plasma glucose with known, switchable
error sources: ISF kinetics (1st order, 2nd order, time-varying tau),
on-device filter tau\_f and pure delay d, calibration drift, night
compression lows, slow plus fast AR(1) noise, sampling at 1/5/15 min and
burst dropouts. Seeds make every trace reproducible.

\begin{longtable}[]{@{}llll@{}}
\toprule\noalign{}
Scenario & ISF kinetics & Extra sensor chain & Artifacts \\
\midrule\noalign{}
\endhead
\bottomrule\noalign{}
\endlastfoot
S0 matched & 1st order, tau known & none & none \\
S1 sensor chain & 1st order & tau\_f 3 min + delay 2 min (unmodelled) &
none \\
S2 second order & two stages, same total lag & none & none \\
S3 tau varying & 1st order, tau x exp(0.25 z), 12 h correlation & none &
none \\
S4 calibration/artifact & 1st order & none & gain drift sd 4\%, offset
drift sd 4 mg/dL, compression lows \\
S5 full & S3 & S1 & S4 \\
\end{longtable}

\hypertarget{tests-11-of-11-pass}{%
\subsection{Tests (11 of 11 pass)}\label{tests-11-of-11-pass}}

\begin{itemize}
\tightlist
\item
  Ramp at 2 mg/dL/min, tau 10: steady error 20.02 (S0), 30.07 (S1),
  20.07 (S2) mg/dL; theory 20, 30, 20.
\item
  Own ISF integrator reproduces simglucose\textquotesingle s
  subcutaneous state: RMSE 0.002 to 0.008 mg/dL, max 0.04 mg/dL (3
  patients).
\item
  Noise sd 4.96 for target 5; autocorrelation matches 0.6 exp(-h/30) +
  0.4 exp(-h/2) within 0.01 at 1, 5, 15, 30 min.
\item
  tau varying: sd(log tau) 0.22 (target 0.25, clipped to 3 to 40 min);
  dropout 4.8\% (target 5\%); compression depth within 0.25 to 0.45;
  seeds reproducible; no index shift when tau tends to 0.
\end{itemize}

\hypertarget{scenario-report-raw-signal-noise-sd-5-mgdl-5-min-sampling-5-dropouts-median-over-27-patients}{%
\subsection{Scenario report (raw signal, noise sd 5 mg/dL, 5-min
sampling, 5\% dropouts, median over 27
patients)}\label{scenario-report-raw-signal-noise-sd-5-mgdl-5-min-sampling-5-dropouts-median-over-27-patients}}

\begin{longtable}[]{@{}lllll@{}}
\toprule\noalign{}
Scenario & RMSE (mg/dL) & MARD (\%) & Best-shift lag (min) & Nominal tau
+ tau\_f + d (min) \\
\midrule\noalign{}
\endhead
\bottomrule\noalign{}
\endlastfoot
S0 & 8.7 & 4.9 & 10 & 11.2 \\
S1 & 11.2 & 5.5 & 15 & 16.2 \\
S2 & 8.9 & 5.0 & 11 & 11.2 \\
S3 & 8.8 & 4.7 & 10 & 11.2 \\
S4 & 11.2 & 7.0 & 10 & 11.2 \\
S5 & 14.1 & 7.2 & 15 & 16.2 \\
\end{longtable}

S0 RMSE 8.7 equals the prototype\textquotesingle s raw RMSE at noise
scale 0, as expected. "Best-shift lag" is the integer shift that
minimises RMSE on noise-free data; it sits about 1 min below the nominal
lag because a first-order lag is not a pure delay.

\hypertarget{findings-that-change-the-plan}{%
\subsection{Findings that change the
plan}\label{findings-that-change-the-plan}}

\begin{enumerate}
\def\labelenumi{\arabic{enumi}.}
\tightlist
\item
  Realism: S5 raw MARD is 6.2, 7.3, 8.7 and 10.7\% at noise sd 2, 5, 8,
  12 mg/dL. One clinical accuracy study reports 9.08\% against venous
  reference (AiDEX, PMC9846405, details only partly verified), so use
  noise sd 8 mg/dL as the stress level; 2 and 5 are optimistic. This is
  a loose anchor only.
\item
  Rapid change is rare and uneven: per patient, time with
  \textbar rate\textbar{} \textgreater= 1 mg/dL/min has median 8.9\%
  (range 0.1 to 27.8\%); \textgreater= 2 mg/dL/min median 3.1\%. Some
  virtual patients barely move, so lag is nearly invisible for them.
  Report results stratified by this fraction and use the rapid-change
  subset for H1.
\item
  S1 and S5 add 5 min of unmodelled lag. An estimator that knows only
  tau will be systematically behind by about 5 min there; this is the
  physiology-versus-sensor identifiability test (H3).
\end{enumerate}

\hypertarget{open-decisions}{%
\subsection{Open decisions}\label{open-decisions}}

\begin{itemize}
\tightlist
\item
  Meal sizes: the prototype\textquotesingle s meals give median peak 243
  mg/dL (range 138 to 336); confirm this is acceptable or rescale.
\item
  Noise sd default for headline results (chosen: 8 mg/dL, with 2 and 5
  as sensitivity).
\end{itemize}
\end{document}
